\documentclass[UTF8,a4paper,12pt]{ctexart}

\usepackage[margin=2.35cm]{geometry}
\usepackage{amsmath,amssymb,bm}
\usepackage{booktabs,longtable,tabularx,array,multirow}
\usepackage{xcolor}
\usepackage{enumitem}
\usepackage{listings}
\usepackage{hyperref}
\usepackage{fancyhdr}
\usepackage{float}

\definecolor{reportgreen}{RGB}{36,92,68}
\definecolor{reportgray}{RGB}{245,247,246}
\hypersetup{
  hypertexnames=false,
  colorlinks=true,
  linkcolor=reportgreen,
  urlcolor=reportgreen,
  citecolor=reportgreen,
  pdfauthor={CUMCM 备赛小组},
  pdftitle={2020--2025 年 CUMCM A/B/C 题学习与论文写作报告}
}
\pagestyle{fancy}
\fancyhf{}
\fancyhead[L]{CUMCM A/B/C 题备赛学习报告}
\fancyhead[R]{2020--2025}
\fancyfoot[C]{\thepage}
\setlength{\headheight}{15pt}
\setlength{\parindent}{2em}
\setlength{\parskip}{0.35em}
\setlist{nosep,leftmargin=2em}
\lstset{
  basicstyle=\ttfamily\small,
  frame=single,
  breaklines=true,
  columns=fullflexible,
  showstringspaces=false
}

\title{\bfseries 2020--2025 年全国大学生数学建模竞赛\\A/B/C 题学习与论文写作报告}
\author{队长、组员1、组员2}
\date{2026 年 7 月 28 日}

\begin{document}
\maketitle

\begin{abstract}
本报告面向三人本科生队伍的全国大学生数学建模竞赛备赛。研究范围严格限定为
2020--2025 年 A、B、C 三类题，论文语料为工作区中 59 篇编号 PDF，共 2892 页；
D、E 题不进入统计。报告先按机理建模、综合优化和数据决策三条主线划分个人责任，
再把思维导图中的算法逐项落实为原理、步骤、代码接口、适用边界和验收产物。论文
分析将编号论文、赛题原文和专家评述分为三类证据：编号论文用于结构和文风统计，
赛题原文用于核对任务，专家评述只用于补充题型背景。全量页面采用文本层提取、
RapidOCR 和 Tesseract 复识相结合的流程，并对公式、表格和小字号图注回看原始栅格。
在此基础上形成十四类章节写作规范、A/B/C 题微观行文差异、创新策略、失分风险与
四周量化训练安排，同时记录双模式论文模板和写作审计 Skill 的实际完成路径。

\textbf{关键词：}数学建模竞赛；机理建模；组合优化；数据决策；论文写作；证据台账
\end{abstract}

\tableofcontents
\newpage

\section{研究边界、证据与处理方法}

\subsection{语料边界}

本报告把名称匹配 A/B/C 加编号的 PDF 视为用户提供的国赛一等奖论文语料，不再
外推奖项名单。按年份和题型清点结果见表~\ref{tab:corpus-count}。59 篇论文全部位于
各年份目录根部，页数合计 2892。原始目录 \texttt{2020/} 至 \texttt{2025/} 及思维
导图 \texttt{1.png} 只读使用，不因 OCR、统计或排版而改写。

\begin{table}[htbp]
\centering
\caption{编号论文语料分布}
\label{tab:corpus-count}
\begin{tabular}{crrrr}
\toprule
年份 & A 题 & B 题 & C 题 & 合计\\
\midrule
2020 & 4 & 4 & 5 & 13\\
2021 & 3 & 4 & 4 & 11\\
2022 & 3 & 3 & 2 & 8\\
2023 & 4 & 3 & 3 & 10\\
2024 & 5 & 3 & 4 & 12\\
2025 & 1 & 2 & 2 & 5\\
\midrule
合计 & 20 & 19 & 20 & 59\\
\bottomrule
\end{tabular}
\end{table}

三类材料的用途不可互换。编号论文回答“获奖论文怎样组织和表达”；赛题原文回答
“当年要求解决什么”；专家评述帮助理解命题背景。后两类材料不进入获奖论文的
章节比例、段落密度、图表与公式安排或高频表述统计。2020 年本地缺少 A/B/C 赛题
原文，表~\ref{tab:problems} 只据专家评述确认题名，不据评述反推子问。

\subsection{页面识别与证据台账}

处理流程以论文为单位建立 JSONL 页面记录。唯一可用文本层的论文直接提取，其余
58 篇先按 120 DPI 渲染并由 RapidOCR 全量识别；文本密集页出现 OCR 中位置信度
低于 0.75、中文正文少于 100 字或标题段落明显断裂时，按 220 DPI 重渲染，并以
Tesseract 的 \texttt{chi\_sim+eng} 复识。每页记录年份、题型、论文编号、页码、
识别方法、置信度、中文字符数、文本框坐标及复识审计信息。

台账把证据细分为章节命中、模型链、写作特征和页面视觉核验四类。公式、表格及
小字号图注不依据 OCR 猜测；这类内容只通过原始栅格的局部裁剪核对。统计脚本必须
先验证 59 个文件和 2892 个连续页码，再计算章节顺序、页数分布、段落密度、图表
候选、公式候选和模型词链。最终全篇页数中位数为 45 页，四分位范围为
37--58.5 页；每篇论文的页面文本行数中位数再取中位数为 38，段落候选数对应
中位数为 6。段落候选
由正文行缩进和垂直间距估计，只用于版面密度比较，不等同于排版软件中的自然段。
共有 2232 页触发并完成 Tesseract 复识，其中 1583 页因质量分更高而被选为最终
文本；最终方法还包括 1281 页 RapidOCR 和 28 页直接文本提取。按同一严格阈值，
1770 页仍标记为未确认，因此这些页面只提供定位候选，不能替代原始栅格核验。
十四类章节的检出率和页级跨度见表~\ref{tab:section-stats}。

\section{六年真题任务图谱与人员分工}

\subsection{十八道真题的训练定位}

表~\ref{tab:problems} 给出真题与主责能力的对应关系。这里的“主责”表示首先完成
建模和代码基线的人，不表示其他成员不参与复核。每题至少安排一名交叉复核者，
避免模型、数据和论文三条链彼此脱节。

\begin{longtable}{p{1.2cm}p{1.3cm}p{3.3cm}p{6.0cm}p{1.9cm}}
\caption{2020--2025 年 A/B/C 真题训练映射}\label{tab:problems}\\
\toprule
年份 & 题型 & 题目 & 训练任务与模型接口 & 主责\\
\midrule
\endfirsthead
\toprule
年份 & 题型 & 题目 & 训练任务与模型接口 & 主责\\
\midrule
\endhead
2020 & A & 炉温曲线 & 仅确认题名；作为传热机理、温度场数值求解与工艺约束训练入口，不虚构本地缺失的子问。 & 队长\\
2020 & B & 穿越沙漠 & 仅确认题名；作为资源约束、路径与多阶段决策训练入口，不虚构子问。 & 组员1\\
2020 & C & 中小微企业信贷策略 & 仅确认题名；作为风险特征、信用评价与额度决策训练入口，不虚构子问。 & 组员2\\
2021 & A & FAST 主动反射面形状调节 & 几何关系、受力/位移约束、连续曲面离散、拟合误差和可实现调节量。 & 队长\\
2021 & B & 乙醇偶合制备 C4 烯烃 & 实验因素建模、反应条件优化、多目标权衡及最优条件解释。 & 组员1，组员2复核\\
2021 & C & 生产企业原材料订购与运输 & 供应商评价、订购计划、运输损耗和多周期约束的联立决策。 & 组员1\\
2022 & A & 波浪能装置最大输出功率设计 & 动力学方程、阻尼与刚度参数、数值积分及平均功率寻优。 & 队长\\
2022 & B & 无人机纯方位无源定位 & 方位几何、编队位置误差、迭代校正与多机协同策略。 & 队长，组员1复核\\
2022 & C & 古代玻璃制品成分分析 & 缺失与异常处理、类别差异、成分关系和未知样品分类。 & 组员2\\
2023 & A & 定日镜场优化设计 & 光学/几何效率、遮挡与截断、场地约束及连续—离散联合优化。 & 队长，组员1复核\\
2023 & B & 多波束测线问题 & 海底几何、覆盖宽度、重叠率、测线方向和总航程优化。 & 组员1，队长复核\\
2023 & C & 蔬菜类商品定价与补货决策 & 销量关联、需求预测、损耗、成本加成与日级补货定价。 & 组员2，组员1复核\\
2024 & A & 板凳龙闹元宵 & 螺线与转向几何、碰撞约束、速度传播和队形可行性。 & 队长\\
2024 & B & 生产过程决策问题 & 零配件抽检、装配、拆解、退换损失和多阶段策略比较。 & 组员1\\
2024 & C & 农作物种植策略 & 地块—季节—作物约束、不确定价格产量和多年度轮作计划。 & 组员1，组员2复核\\
2025 & A & 烟幕干扰弹投放策略 & 三维运动、遮蔽几何、时空覆盖和多平台协同投放。 & 队长，组员1复核\\
2025 & B & 碳化硅外延层厚度确定 & 光谱信号处理、物理测量模型、参数反演和误差分析。 & 队长，组员2复核\\
2025 & C & 无创产前检测时点与异常判定 & 分组时点优化、异常识别、风险约束和模型泛化检验。 & 组员2\\
\bottomrule
\end{longtable}

\subsection{四块学习结构}

队伍采用“三条主责线加一条交叉线”。队长主责机理建模、常微分/偏微分方程、
物理模型、RK4、有限差分、有限元、模型验证和代码集成；组员1主责规划、整数和
多目标优化、智能优化、图论、网络优化、仿真及灵敏度分析；组员2主责数据清洗、
探索、降维、评价决策、回归分类、时间序列和交叉验证。三人共同学习
MATLAB/Python、量纲与误差分析以及 A/B/C 题型选型。共享不等于无人负责：公共
模块仍由对应主责人提供接口，另两人按统一测试样例复核。

\section{理论知识学习笔记}

\subsection{队长：机理、连续模型与数值实现}

\subsubsection{覆盖内容、原理与求解步骤}

队长的模型链从守恒关系开始，而不是从求解器开始。机械系统用牛顿第二定律、
刚体平衡、重力、摩擦和弹簧—阻尼关系建立状态方程；流体问题用质量守恒和
Bernoulli 关系闭合；电路问题用 Ohm 定律和 Kirchhoff 定律组织节点或回路方程；
气体过程核对理想气体状态关系。传热采用 Newton 冷却或热传导方程，波动、扩散
及多孔介质渗流按偏微分方程处理。人口过程用 Logistic 方程和捕食—被捕食模型，
药物过程使用房室代谢方程，传播过程比较 SIR 与 SIS 的恢复机制。

完整步骤固定为：
\begin{enumerate}
  \item 画出对象、边界和能量/质量/动量交换，确定系统尺度与正方向；
  \item 从守恒律或受力关系推导状态方程，逐项给出单位和符号；
  \item 写明初始条件、边界条件、连续性条件、控制量及可行域；
  \item 先做极端情形、量纲和简化模型的解析检查；
  \item 对 ODE 比较显式 Euler 与四阶 Runge--Kutta，对 PDE 按网格选择有限差分或有限元；
  \item 进行步长/网格收敛、守恒残差、回代误差和参数扰动检验；
  \item 把经检验的状态量封装为组员1的目标/约束或组员2的可解释特征。
\end{enumerate}

Euler 法实现简单但全局误差阶数低，只用于基线和步长敏感性对照；RK4 在光滑 ODE
上兼顾精度与实现成本，但对刚性系统不应硬套。有限差分适合规则区域和清晰的
差分格式，可直接检查截断误差与 CFL 条件；有限元更适合复杂几何和变系数边界，
代价是网格、形函数、装配和边界处理更重。解析解存在时先用解析解作单元测试，
不存在时采用网格加密、守恒量和独立实现交叉验证。

\subsubsection{开源代码逻辑、约束与选型边界}

代码流水线按下列接口拆分：
\begin{quote}
\small\ttfamily
geometry -> state -> rhs -> solver\\
-> validation -> export
\end{quote}
MATLAB 侧以 \texttt{ode45}、稀疏矩阵和优化工具箱形成对照实现；Python 侧
以 NumPy/SciPy 的数组、\texttt{solve\_ivp}、稀疏线性代数和网格数据结构实现。
每个模型提供一个可手算的小样例和一个真实规模样例，输出时间网格、状态量、
守恒残差、收敛曲线及下游所需的 CSV。

模型约束至少包括变量域、初边值、守恒关系、几何/运动学限制、参数物理范围和
数值稳定条件。能由规则直接计算的几何关系不使用神经网络；样本驱动关系没有
物理闭合时，不伪造微分方程。ODE 与 PDE 的边界取决于空间耦合是否不可忽略；
有限差分与有限元的边界取决于几何和介质复杂度；RK4 与自适应/刚性求解器的
边界取决于稳定域和尺度差异。

\paragraph{可实现代码清单}
\begin{itemize}
  \item 牛顿运动、弹簧阻尼、刚体平衡、摩擦与重力模块；
  \item Logistic、捕食—被捕食、SIR/SIS、药物房室、Newton 冷却模块；
  \item 热、波、扩散和多孔流方程的 FDM 基线及 FEM 对照；
  \item Euler、RK4、自适应积分接口和统一误差指标；
  \item 网格加密、守恒残差、量纲检查、参数回代和结果导出；
  \item 2022A、2023A、2024A、2025A/B 的最小复现实验。
\end{itemize}

\subsection{组员1：规划、智能优化、图网络与仿真}

\subsubsection{覆盖内容、原理与求解步骤}

规划模型先按变量和约束结构分类：连续线性规划、非线性规划、二次规划、整数
规划、混合整数线性规划和 0--1 规划；再区分单目标与多目标、确定性与随机性。
线性或凸结构优先使用有最优性证书的确定性求解器。整数变量表示启停、选择、
顺序和分配时，必须写出 Big-M 的来源或改用指示约束，不能以取整替代整数规划。
多目标问题先报告 Pareto 关系，再依据题意选择加权、分层或 $\varepsilon$ 约束，
不凭个人偏好给权重。

智能优化覆盖遗传算法、粒子群、模拟退火、蚁群、差分进化和贪心基线，以及
NSGA-II、MOPSO。思维导图中的鲸鱼优化、麻雀搜索、蛾火焰优化作为受控对照：
只有在变量编码、可行性修复和公平预算明确后才进入实验。复杂名称不构成创新，
必须与确定性基线或简单启发式在同一目标、随机种子组和函数评估次数下比较。

图论部分覆盖 Dijkstra、Floyd、SPFA，Prim、Kruskal，最大流、最小费用流、
路径规划、二分图匹配和节点重要性。稀疏非负边单源问题优先 Dijkstra；全源且
规模中等用 Floyd；存在负边时先检查负环，SPFA 只作工程候选而非复杂度保证。
最小生成树解决连通成本，不替代最短路；最大流回答容量，最小费用流同时回答
流量和成本；指派关系优先二分图匹配。

统一求解步骤为：
\begin{enumerate}
  \item 定义决策变量、目标、硬约束、软约束和不确定参数；
  \item 建立可行性基线，并在小规模枚举样例上核对目标值；
  \item 判断线性、凸性、整数性、网络结构和多目标冲突；
  \item 先运行确定性/贪心基线，再决定是否启用元启发式；
  \item 写清编码、初始化、选择更新、约束修复、终止条件和随机种子；
  \item 记录可行率、最优界/间隙、收敛轨迹、运行时间和跨种子分布；
  \item 通过 Monte Carlo 或情景仿真传播参数不确定性，并做敏感性分析。
\end{enumerate}

\subsubsection{开源代码逻辑、优缺点与边界}

代码流水线为：
\begin{quote}
\small\ttfamily
data -> instance -> model -> solver -> repair\\
-> evaluate -> simulate -> report
\end{quote}
Python 用 PuLP/Pyomo 或 SciPy 建立确定性基线，用
NetworkX 验证图算法接口；MATLAB 用 Optimization Toolbox/Global Optimization
Toolbox 建立第二实现。随机算法统一接收种子、预算和停止阈值，输出每代最优值、
可行率及 Pareto 集。仿真模块只通过明确的数据类接收方案，禁止在仿真循环内
偷偷改变约束。

精确规划的优势是可解释、可复算并可给上下界，缺点是非凸或大规模整数问题可能
耗时；启发式的优势是能处理混合和非光滑结构，缺点是没有全局最优保证且对参数
敏感。模型边界由结构而非题型标签决定：可线性化的约束先线性化；不确定性已知
分布可用随机规划或 Monte Carlo，只有区间信息时优先稳健情景；动态图路径与
静态最短路分开处理。

\paragraph{可实现代码清单}
\begin{itemize}
  \item LP、NLP、QP、整数/MILP、0--1 模型生成器及小规模枚举核验；
  \item GA、PSO、SA、ACO、DE、贪心和 NSGA-II/MOPSO 公平比较框架；
  \item WOA、SSA、MFO 的受控插件与相同函数评估预算；
  \item Dijkstra、Floyd、SPFA、Prim、Kruskal、最大流、最小费用流和二分图匹配测试；
  \item Monte Carlo/离散事件仿真、单因素与全局敏感性接口；
  \item 2020B、2021C、2023B、2024B/C 的可行基线。
\end{itemize}

\subsection{组员2：数据处理、评价决策与统计学习}

\subsubsection{数据链、算法原理与完整步骤}

数据工作从数据字典和泄漏边界开始。缺失值按形成机制选择均值/中位数、插值、
KNN 填补或删除；异常值用 $3\sigma$、箱线图和 Isolation Forest 交叉检查，
不能把业务极值自动删除；重复记录先定义主键和时间容差。尺度处理比较
Z-score 标准化与 Min--Max 归一化，类别变量按名义/有序属性选择 one-hot 或
label 编码，连续变量只有在解释或规则需要时才按等宽/等频分箱。

探索阶段采用 Pearson、Spearman、Kendall 相关，分别对应线性、单调秩和序关系；
聚类比较 K-Means、DBSCAN、层次聚类和谱聚类。空间或时间缺口比较 Lagrange、
三次样条和 Kriging 插值。PCA 用于线性方差压缩，t-SNE 主要用于局部结构可视化，
UMAP 兼顾局部与部分全局结构；后两者的二维图不直接充当预测证据。

评价决策链覆盖 AHP、模糊综合评价、模糊 AHP、灰色关联、PCA、因子分析、熵权、
CRITIC、变异系数、TOPSIS、VIKOR、灰色评价和物元可拓，并训练熵权--AHP、
CRITIC--TOPSIS 等组合。主观权重必须给判断矩阵和一致性检验，客观权重必须解释
数据离散或冲突如何影响权重；TOPSIS 的距离、VIKOR 的折中和灰色关联的序列相似
并非同一含义，不能只因都输出排序而互换。

回归覆盖多元线性、Ridge、Lasso、SVR、随机森林、XGBoost 和 CatBoost；分类覆盖
Logistic、朴素 Bayes、SVM、决策树、随机森林、AdaBoost 和 CNN。时间序列覆盖
ARIMA、SARIMA、指数平滑、VAR、GM(1,1)、灰色--Markov、Prophet、XGBoost、
LightGBM、BP、RNN、LSTM、TCN 和 GRU。线性/广义线性模型保留可解释基线，树模型
处理非线性与交互，深度网络仅在样本规模、序列长度和验证方案足以支撑时启用。

完整步骤为：
\begin{enumerate}
  \item 建立字段、单位、主键、时间戳、缺失和异常审计表；
  \item 在任何填补、标准化、降维或特征选择前划分训练/验证边界；
  \item 完成探索图表与业务一致性核验，建立简单均值/线性/持续性基线；
  \item 按任务选择评价、回归、分类、聚类或时序模型，并列出近邻方法边界；
  \item 将预处理与模型封装在同一 pipeline，采用 K 折、留出或滚动窗口验证；
  \item 回归报告 MAE、MSE、RMSE、MAPE，分类报告混淆矩阵及类别不平衡指标；
  \item 做残差、校准、特征稳定性、阈值和数据扰动分析，再向决策模块输出。
\end{enumerate}

\subsubsection{开源代码逻辑、约束与选型边界}

代码流水线为：
\begin{quote}
\small\ttfamily
schema -> audit -> split -> preprocess -> baseline\\
-> fit -> validate -> explain -> export
\end{quote}
Python 侧使用 pandas、scikit-learn、
statsmodels 及主流梯度提升接口；MATLAB 侧建立读表、缺失处理、回归分类和时序
对照。所有随机切分记录种子，时间序列禁止随机打乱，分组样本按主体隔离。
评价模型在正向化、无量纲化和权重确定后再计算排序，并做权重扰动和秩相关检验。

K-Means 假定近似球形且需预设簇数，DBSCAN 能识别噪声和任意形状但依赖密度尺度；
层次聚类便于解释树状结构但大样本代价高；谱聚类适合非凸簇但需构图。Lasso 可
稀疏选择，Ridge 对共线性更稳；SVR 适合中小样本但调参和扩展性受限；随机森林
稳健易用但外推弱；梯度提升精度高但要防泄漏和过拟合。ARIMA/SARIMA 要检验
平稳和季节性，VAR 要控制维数和滞后，GM(1,1) 只适合小样本单调趋势，深度时序
模型不能替代滚动窗口基线。

\paragraph{可实现代码清单}
\begin{itemize}
  \item 缺失、异常、重复、尺度、编码、分箱和数据泄漏审计 pipeline；
  \item Pearson/Spearman/Kendall、四类聚类、三类插值和 PCA/t-SNE/UMAP；
  \item AHP/FAHP、熵权、CRITIC、灰色关联、TOPSIS、VIKOR、物元可拓和组合评价；
  \item 线性/Ridge/Lasso、SVR、RF、XGBoost/CatBoost 回归及分类基线；
  \item ARIMA/SARIMA、指数平滑、VAR、GM/灰色--Markov、Prophet 和神经时序对照；
  \item K 折、留出、分组、滚动窗口验证与统一指标报告；
  \item 2020C、2022C、2023C、2025C 的数据基线。
\end{itemize}

\subsection{三人交叉学习与代码集成}

三人共同维护数据契约、配置文件、日志和结果目录。MATLAB 适合快速矩阵推导、
数值求解与官方工具箱复核；Python 适合数据 pipeline、开源模型、自动化测试和
制图。关键结果至少满足“同一语言的独立小样例”或“MATLAB/Python 双实现”之一。
误差分析统一区分测量误差、参数误差、模型结构误差、离散误差和随机算法波动，
不能只给一个综合误差率。

选题时先在两小时内完成任务矩阵：状态量和守恒关系占主导时优先 A 类机理线；
离散资源、路径、排程和策略占主导时优先 B 类优化线；附件数据、预测、评价和
不确定决策占主导时优先 C 类数据线。最终决定同时考虑题意闭合度、可验证性、
团队代码储备和结果文件要求，而不是依据“哪题看起来熟悉”。

\subsection{阶段成果、真实短板与四周计划}

\subsubsection{当前可取之处}

本轮已经形成全算法责任表、18 题训练接口、统一代码分层、模型选型边界和论文
证据流水线。可取之处不是“算法数量多”，而是每条主责线都要求基线、约束、
独立检验和下游接口；复杂模型必须与简单模型公平比较；模型、代码和论文中的
变量名与结果文件可以相互追溯。

\subsubsection{当前真实短板}

成员未提供个人既有学习记录，因而本报告不虚构“已经熟练掌握”的经历。当前能
确认的是计划与分析产物，不能确认的是三人在限时条件下的独立实现能力。具体
风险包括：队长尚缺复杂边界 PDE/FEM 的限时网格调试证据；组员1尚缺大规模
MILP 的界、间隙与不可行诊断记录，也缺元启发式跨随机种子统计；组员2尚缺严格
时间/分组隔离的数据验证记录，深度时序和客观评价权重容易停留在调用层面。
全队尚未完成一场 72 小时全流程演练，论文数字、图表和附件代码的一致性也没有
实战验收。这些短板按“尚无证据”表述，不把未知情况写成个人缺陷。

\subsubsection{2026 年 7 月 29 日至 8 月 25 日量化训练}

\begin{longtable}{p{1.4cm}p{1.55cm}p{3.1cm}p{4.6cm}p{3.4cm}}
\caption{四周训练计划与验收标准}\label{tab:four-weeks}\\
\toprule
时间 & 负责人 & 训练任务 & 产物 & 验收标准\\
\midrule
\endfirsthead
\toprule
时间 & 负责人 & 训练任务 & 产物 & 验收标准\\
\midrule
\endhead
7/29--8/4 & 队长 & 完成 2022A 动力学最小复现；比较 Euler、RK4 与自适应积分。 & 方程说明、两套代码、步长收敛图、守恒与回代误差表。 & 小样例误差随步长按预期下降；关键结果可复算；接口测试全部通过。\\
7/29--8/4 & 组员1 & 完成 2024B 的枚举基线与 0--1/MILP 表达。 & 决策树、模型文件、枚举核验、不可行情形清单。 & 小规模最优值一致；每条约束有题意来源；报告最优界或间隙。\\
7/29--8/4 & 组员2 & 完成 2022C 数据审计、可视化与分类基线。 & 数据字典、清洗审计、pipeline、验证报告。 & 预处理只在训练集拟合；异常处理可追溯；至少三类基线同口径比较。\\
8/5--8/11 & 队长 & 完成热/扩散方程 FDM 与复杂几何 FEM 对照。 & 网格文件、边界测试、网格收敛表。 & 至少三档网格；边界和量纲单测通过；说明 FDM/FEM 选型边界。\\
8/5--8/11 & 组员1 & 完成 2023B 图/几何基线及 GA、PSO 公平对照。 & 基线求解器、随机算法报告、可行性修复测试。 & 不少于 20 个种子；预算相同；报告中位数、离散度、时间和可行率。\\
8/5--8/11 & 组员2 & 完成 2023C 时序滚动验证和补货特征接口。 & 持续性/季节基线、ARIMA、树模型、滚动评估。 & 无未来信息泄漏；至少四个误差指标；失败日期有原因分析。\\
8/12--8/18 & 全队 & 选择一题进行 36 小时压缩演练，贯通模型、代码、图表、论文和附件。 & 可编译论文、运行入口、结果清单、复盘表。 & 所有正文数字可定位到输出；无未定义引用；第三人按说明复现关键表。\\
8/19--8/25 & 全队 & 进行 72 小时 A/B/C 三选一全流程模拟并交叉审稿。 & 最终论文、源代码、附件、审计 JSON、个人复盘。 & 12 小时内定题；24 小时出基线；48 小时主结果冻结；终稿通过结构、引用、检验和复现门。\\
\bottomrule
\end{longtable}

\section{59 篇编号论文的结构与写作分析}

\subsection{宏观框架与版面组织}

编号论文的共同主线不是固定的小标题数量，而是从任务到证据的闭环：
\[
\begin{aligned}
\text{任务界定}&\longrightarrow\text{假设与变量}
\longrightarrow\text{模型及约束}\longrightarrow\text{求解实现},\\
&\longrightarrow\text{结果解释}
\longrightarrow\text{独立检验}\longrightarrow\text{边界与复现材料}.
\end{aligned}
\]
章节标题可以因题目合并，但上述功能不能被删去。篇幅判断使用全量论文的实测
页级位置和章节跨度，不把经验百分比写成硬规则。全篇总页数参考中位数
45 页和四分位范围 37--58.5 页；不同章节的检出率、首次出现位置和跨度汇总于
表~\ref{tab:section-stats}。检测率低的章节只能形成
定性建议，不能据缺失 OCR 推断获奖论文没有相应内容。

\begin{longtable}{p{1.8cm}p{2.5cm}p{3.0cm}p{2.8cm}p{3.0cm}}
\caption{十四类章节的全语料检出与页级位置统计}\label{tab:section-stats}\\
\toprule
章节 & 检出论文 & 首次位置比例中位数 [Q1,Q3] & 跨度页数中位数 [Q1,Q3] & 跨度比例中位数 [Q1,Q3]\\
\midrule
\endfirsthead
\toprule
章节 & 检出论文 & 首次位置比例中位数 [Q1,Q3] & 跨度页数中位数 [Q1,Q3] & 跨度比例中位数 [Q1,Q3]\\
\midrule
\endhead
摘要 & 58/59 (98.3\%) & 0.0215 [0.017, 0.0268] & 1 [1, 1] & 0.0215 [0.017, 0.0268]\\
问题重述 & 57/59 (96.6\%) & 0.0444 [0.0345, 0.0541] & 1 [1, 1] & 0.0222 [0.0172, 0.0286]\\
问题分析 & 57/59 (96.6\%) & 0.0526 [0.0417, 0.0732] & 1 [1, 1] & 0.0244 [0.0172, 0.0357]\\
模型假设 & 48/59 (81.4\%) & 0.0735 [0.0562, 0.106] & 1 [1, 1] & 0.0225 [0.0172, 0.0286]\\
符号说明 & 58/59 (98.3\%) & 0.0808 [0.0661, 0.1143] & 1 [1, 3.75] & 0.029 [0.0205, 0.0749]\\
模型建立 & 58/59 (98.3\%) & 0.0825 [0.0489, 0.1329] & 2 [1, 6] & 0.045 [0.0245, 0.1681]\\
模型求解 & 55/59 (93.2\%) & 0.1077 [0.0286, 0.2319] & 2 [1, 6] & 0.0385 [0.0206, 0.1562]\\
结果分析 & 47/59 (79.7\%) & 0.1795 [0.0639, 0.3107] & 7 [1, 12] & 0.1282 [0.0278, 0.25]\\
模型检验 & 22/59 (37.3\%) & 0.2294 [0.0564, 0.4821] & 1.5 [1, 7] & 0.0335 [0.0249, 0.1889]\\
灵敏度分析 & 25/59 (42.4\%) & 0.0357 [0.02, 0.3929] & 1 [1, 2] & 0.0256 [0.0179, 0.0385]\\
模型评价 & 51/59 (86.4\%) & 0.5111 [0.396, 0.6186] & 1 [1, 1] & 0.0227 [0.0172, 0.0294]\\
改进方案 & 19/59 (32.2\%) & 0.4667 [0.3707, 0.5928] & 1 [1, 1] & 0.0263 [0.0194, 0.0345]\\
参考文献 & 57/59 (96.6\%) & 0.5424 [0.4151, 0.6667] & 1 [1, 11] & 0.0333 [0.0222, 0.2941]\\
附录 & 48/59 (81.4\%) & 0.5574 [0.4427, 0.7079] & 16 [6, 28] & 0.3779 [0.1865, 0.4926]\\
\bottomrule
\end{longtable}

表中首次位置按页码除以论文总页数计算。灵敏度分析的首次位置中位数异常靠前，
抽查发现目录页和摘要中的同名词会造成提前命中，故该项只用于检索，不能解释为
正文真的在前 3.57\% 篇幅内展开灵敏度分析。参考文献与附录跨度也可能因标题合并、
附件页连续而偏大；模板只把高检出章节的跨度用于量化提示。

段落组织以“一个判断配一组证据”为单位。建模段通常先说明关系来源，再定义变量
并列方程，公式后解释每一项的作用；求解段按输入、初始化、迭代、停止和输出展开；
结果段按读数、比较、原因和决策含义展开。图表应在首次讨论附近出现，正文必须
引用并解释，完整中间矩阵和长代码移至附录。公式、表格和图注的数量只能作为
页面候选统计，最终含义以原始页面裁剪核验为准。每篇图表标题候选数的中位数及
四分位范围为 17 [14, 23]，公式行候选对应为 420 [228.5, 511.5]；两类候选页在
全文相对位置的中位数分别为 29.8\% 和 56.9\%。这些计数用于判断版面节奏，不表示
精确的图表或公式数量。

\subsection{A/B/C 三类论文的微观文风}

分题型比较采用各类论文内部的章节、连接词和模型组论文检出数，不将某个模型词
出现次数直接解释为模型质量。A 类总页数中位数和四分位范围为 44.5 [36.5, 58]
页，模型组检出前三项为机理与数值 15 篇、规划与智能优化 13 篇、时序预测 3 篇；
B 类为 42 [34, 55] 页，前三项为机理与数值 12 篇、规划与智能优化 10 篇、统计
学习 6 篇；C 类为 52 [43.25, 60.25] 页，前三项为时序预测 17 篇、规划与智能
优化 17 篇、统计学习 13 篇。模型词可能来自摘要、比较方法或参考文献，这里只报告
“论文中检出”，不把它直接解释为论文主模型。

\subsubsection{A 类：先交代机制，再讨论数值}

A 类行文从研究对象、坐标、尺度和守恒关系进入。高信息密度表达包括“以……为
研究对象，取……方向为正”“由质量/能量/动量守恒可得”“在边界……处满足……”
“将区域离散为……，当网格由……加密至……时，目标量变化……”。这里的关键
不是多写推导，而是使假设、方程、初边值和数值格式彼此闭合。结果解释要回到
物理量、数量级和机制，不把一条平滑曲线等同于模型正确。

需要避开的写法是先报“使用微分方程模型”，随后直接给软件输出；或者给出方程
却没有边界、单位、稳定条件和网格收敛。图形以结构示意、网格、状态曲线和残差/
收敛图为主，颜色只服务于变量区分。

\subsubsection{B 类：变量、约束与策略必须逐项对应}

B 类先写决策对象、目标口径和约束来源。典型信息链是“令……表示是否选择……”
“由产能/路径/时序要求得到约束……”“为比较冲突目标，构造……并保留 Pareto
方案”“在相同函数评估预算下，改进方法相对基线……”。每个二元变量应能翻译回
一个实际动作，每个约束都能回到题意，每个输出都能转化为方案表。

不能只给元启发式流程图和最优值。论文需报告编码、初始化、约束修复、参数、
终止条件、随机种子或多次运行统计；能使用精确规划时，还应报告界或间隙。图表
以决策结构、算法流程、收敛/Pareto 图、方案对照和敏感性图为主。

\subsubsection{C 类：数据口径、隔离验证与决策解释并重}

C 类从字段、样本单位、时间范围、缺失异常和训练/验证边界写起。有效表达包括
“以……为一条观测，按主体/时间划分训练与验证集”“所有标准化和特征选择仅在
训练折拟合”“相对持续性/线性基线，模型在……指标上……”“该变化主要对应……
特征，但结论限于……范围”。评价排序需解释指标正向化、权重来源和排序稳定性；
预测需把误差转换为补货、时点或风险决策的影响。

需要避开的是全数据清洗后再切分、用随机切分检验未来预测、把相关性写成因果、
只报告最优模型以及用二维降维图证明分类性能。图表以数据质量、分布与相关、
验证协议、残差/校准、特征影响和决策结果为主。

\subsection{十四类章节的写法}

\begin{longtable}{p{1.8cm}p{4.2cm}p{4.0cm}p{4.0cm}}
\caption{章节功能、段落套路与质量门}\label{tab:sections}\\
\toprule
章节 & 必答问题与段落结构 & 推荐表达及图表 & 质量门与常见错误\\
\midrule
\endfirsthead
\toprule
章节 & 必答问题与段落结构 & 推荐表达及图表 & 质量门与常见错误\\
\midrule
\endhead
摘要 & 逐问写对象、模型、关键结果和检验；最后给关键词。 & “针对……，建立……；求得……；经……检验……”；通常不放图。 & 每问均出现，数字与正文一致；不堆算法名，不写“效果良好”。\\
问题重述 & 重组给定量、输出、硬约束、附件和问间依赖。 & “给定……，要求在……约束下确定……”；复杂几何可重绘示意图。 & 不大段照抄，不漏单位、时段或提交文件，不提前引入模型。\\
问题分析 & 每问说明核心矛盾、输入输出、基线、改进和数据流。 & “由于……，故先……，再……”；节末可放总流程图。 & 必须解释选择理由和近邻方法边界，不能只列算法。\\
模型假设 & 一条一义写对象、范围、简化理由和方程后果。 & “在……尺度内忽略……，因此……”；必要时放情景表。 & 题面事实不是假设；不得写“数据真实可靠”；后文可追溯。\\
符号说明 & 定义状态量、决策量、参数、集合、索引、单位和域。 & 三线表分组；局部临时量在首次出现处定义。 & 同符号不跨问改义；公式、代码和附件字段一致。\\
模型建立 & 依次给依据、变量、数学表达、约束、初边值和参数。 & “令……；由……可得……”；结构图靠近定义。 & 公式前有来源、后有解释；约束方向、量纲和变量域正确。\\
模型求解 & 写预处理、初始化、求解器/迭代、参数、停止和输出。 & 按执行顺序写；伪代码和流程图不替代参数表。 & 第三方可复算；不以“调用 MATLAB/Python”结束。\\
结果分析 & 读关键值，比较基线，解释原因，说明约束和决策含义。 & “在……条件下，……由……变为……，表明……”；图表后立即解释。 & 不逐格念表；有单位、有效数字和不确定范围；每问均回应。\\
模型检验 & 明确检验对象、方法、阈值、独立数据、指标和裁决。 & 回代、守恒、残差、网格收敛、交叉验证或消融。 & 不能用训练拟合代替外部能力；失败项也要报告。\\
灵敏度分析 & 定义扰动参数、范围、采样、重复求解和响应指标。 & “当……位于……时结论不变；超过……后策略切换”；放曲线/情景表。 & 扰动后必须重算；区分数值稳定和决策稳定。\\
模型评价 & 将优势绑定机制，将不足绑定假设、数据或算法。 & “由于显式保留……，模型能够……；受……限制……” & 不写任何模型都适用的“简单准确、适用面广”。\\
改进方案 & 对每项局限给新数据/变量/约束、实现方式和验证协议。 & 按“缺陷—改法—验证”对应；可用短表。 & 增加复杂度不是天然改进；方案必须可执行和可检验。\\
参考文献 & 只列正文实际使用的理论、数据、软件与标准来源。 & 统一著录；定义和算法优先原始论文、教材、标准或官方资料。 & 正文引用与条目双向对应；网页含机构、日期和访问日期。\\
附录 & 给入口脚本、环境、参数、随机种子、数据字典和补充结果。 & 文件清单、目录树、可复制代码；主文明确引用。 & 关键结果可由附录复现；不放截图代码，不泄露个人信息。\\
\bottomrule
\end{longtable}

\subsection{可复用表达与“避雷”边界}

可复用的是信息结构，不是整句套用。摘要可用“针对对象和任务—建立核心模型—
给出关键结果—报告检验”的四步；模型段可用“关系依据—变量定义—方程/约束—
式后解释”；结果段可用“读数—比较—机制—决策”；评价段可用“机制带来的能力—
假设造成的边界—可验证改进”。“首先、其次、最后”“通过软件求解可得”“结果
较为理想”“模型具有较强普适性”等句子若没有对象、数值、条件和证据，应删除
或改写。

表达复用不等于隐藏不确定性。OCR 无法确认的字词标为未确认；跨论文统计只报告
检测定义和样本数；本地缺少的 2020 赛题原文不由评述补写。论文中的个别错误也
不应遮蔽其优秀部分：分析重点是任务闭合、关键推导、结果解释和验证等产生高
判别力的做法，同时把错误作为质量门输入。

\section{创新切入、高频失分与高分策略}

\subsection{从基线到差异化模型}

创新先回答“基础方案在哪个可验证环节失败”，再引入复杂度。可执行路径包括：
\begin{enumerate}
  \item \textbf{结构创新：}把题意中的时空、网络、分层或阶段结构显式写入状态和约束；
  \item \textbf{机理—数据融合：}用守恒方程约束参数范围，以数据校准未知项，并分开报告拟合和外推；
  \item \textbf{基线—改进链：}先给解析、贪心、线性或持续性基线，再针对残差、非线性或不确定性改进；
  \item \textbf{约束驱动：}将现实可行性写入模型，并报告活跃约束、不可行诊断和修复代价；
  \item \textbf{不确定性驱动：}把参数误差传播到结果和决策，寻找稳定区间与策略切换点；
  \item \textbf{验证创新：}增加独立回代、网格收敛、消融、滚动验证或反事实情景，使结论比模型名更可信。
\end{enumerate}

不采用“换一个更新颖算法名称”作为独立创新。若鲸鱼、麻雀或蛾火焰优化相对
确定性基线、GA/PSO 没有同预算优势，就不进入主模型；若深度网络没有超过季节
或树模型基线，保留简单方案；若多指标组合权重没有通过扰动和秩稳定检验，不把
组合本身写成贡献。

\subsection{高频失分点}

\begin{enumerate}
  \item 题问、模型、结果表和摘要无法逐一对应，或者漏交附件结果；
  \item 假设冗余、互相冲突，或把题面事实和“数据可靠”写成假设；
  \item 符号未定义、单位不统一、公式逻辑或约束方向错误；
  \item 只介绍算法，不写变量、参数、初边值、约束和停止条件；
  \item 启发式算法只展示一次最好结果，没有基线、预算和随机性统计；
  \item 数据预处理泄漏到验证集，时间序列随机切分，相关关系被写成因果；
  \item 图表缺标题、单位、标签或正文引用，结果段只罗列数字；
  \item 把训练拟合、单点扰动或一条收敛曲线当作完整模型检验；
  \item 优缺点是套话，改进方案只有“增加数据、优化算法”；
  \item 参考文献与正文不对应，附录缺入口、依赖、随机种子或输出说明。
\end{enumerate}

\subsection{高分写作策略与冻结节奏}

高分策略可落实为五个冻结点：定题后冻结任务矩阵；基线通过手算/小规模枚举后
冻结变量和口径；主模型通过独立检验后冻结关键数字；图表与附件可追溯后冻结
结果；最后才冻结摘要。任何改动关键结果的代码都必须重新生成正文表图，避免
人工复制导致不一致。评阅者首先需要看见的是问题闭合、模型合理、数字可核验和
结果有解释，而非模型数量。

提交前执行双人交叉核对：一人只看题问—摘要—结果是否对应，另一人只看假设—
变量—公式—代码—附件是否对应；最后由未撰写该段的人检查引用、单位、有效
数字、图表位置和页面溢出。审计脚本用于发现结构和引用问题，但不能替代公式
正确性、学术判断和原始页面核验。

\section{三项工具成果及迭代机制}

\subsection{学习报告}

本文件 \texttt{assignment.tex} 使用 \texttt{ctexart} 编写，含完整导言区、
标题、正文、公式、长表、参考材料和附录，可由 XeLaTeX 独立编译。报告只陈述
模板和 Skill 的完成情况，不把两者全文嵌入，三项成果保持独立。

\subsection{双模式论文模板}

模板实际路径为 \path{thesis template/}。\par
目录保留上游
\path{cumcmthesis.cls}、字体和
\path{example.tex}；\path{main.tex} 默认为指导模式，在每级标题后给出目标、
清单、语言、篇幅排版、图表位置、判别力要点和错误警示；\texttt{submission.tex}
定义模式开关后载入同一主文件，隐藏全部指引。章节顺序固定为摘要与关键词、
问题重述、问题分析、模型假设、符号说明、分问题模型建立与求解、结果分析与
模型检验、灵敏度/稳健性、模型评价与改进、参考文献、附录。

篇幅提示集中使用本轮 59 篇论文的实测中位数和四分位范围，不预造比例。目录
\texttt{ITERATION.md} 记录证据范围和改造日期；后续范文结论先在批注区登记
来源编号、日期、页码、题型和核验状态，再增量修改指引。

\subsection{MCM-CUP-STANDARD-WRITE Skill}

Skill 位于
\texttt{C:/Users/Lenovo/.codex/skills/mcm-cup-standard-write/}，调用名为
\texttt{\$mcm-cup-standard-write}，界面名保留
\texttt{MCM-CUP-STANDARD-WRITE}。核心文件写入任务规定的两条第 0 条及第 1--7
条约束，详细资料拆为论文结构、A/B/C 题型、章节句式、模型创新、质量门和证据
更新六份 reference。公开审计接口为：
\begin{quote}
\small\ttfamily
python audit\_manuscript.py MAIN.TEX\\
\hspace*{1em}--problem-type A\textbar B\textbar C\\
\hspace*{1em}--format text\textbar json
\end{quote}
脚本检查章节缺失和顺序、占位符、图表标签、交叉引用、文献引用、结果解释、
模型检验及附录结构。它的输出是质量门线索，不能证明模型结论正确。

\section{结论}

本轮工作把备赛从“算法清单”转为可验收的模型、代码和写作链。三名成员分别承担
机理连续模型、综合优化和数据决策的首责，同时共享编程、误差分析和选题判断；
18 道真题为每条能力提供了训练入口。论文研究限定在 59 篇编号语料，统计与页面
证据可追溯，赛题和评述不混入文风样本。模板和 Skill 已形成独立工具，四周计划
把尚无证据的能力转化为代码、图表、论文和复现验收。下一阶段的重点不是继续
增加算法名称，而是在限时演练中证明基线、约束、检验和表达能够稳定闭环。

\begin{thebibliography}{99}
\bibitem{papers}
2020--2025 年全国大学生数学建模竞赛 A/B/C 题编号论文语料（本地归档），
59 篇，2892 页。

\bibitem{problems}
全国大学生数学建模竞赛赛题原文：2021--2025 年 A/B/C 题（本地归档）。

\bibitem{commentaries}
全国大学生数学建模竞赛专家评述：2020--2025 年相关题目（本地归档）。

\bibitem{mindmap}
数学建模学习思维导图 \texttt{1.png}（用户提供，本地文件，核验日期：
2026-07-28）。
\end{thebibliography}

\appendix
\section{算法责任与验收总表}

\begin{longtable}{p{2.1cm}p{5.0cm}p{2.0cm}p{5.5cm}}
\caption{思维导图算法的完整分配}\label{tab:algorithm-owner}\\
\toprule
模块 & 算法或模型 & 主责 & 最小验收\\
\midrule
\endfirsthead
\toprule
模块 & 算法或模型 & 主责 & 最小验收\\
\midrule
\endhead
缺失与异常 & 均值/中位数、插值、KNN、删除；$3\sigma$、箱线图、Isolation Forest；重复样本删除/合并 & 组员2 & 处理前后审计表、规则理由、下游敏感性\\
尺度与编码 & Z-score、Min--Max、one-hot、label、等宽/等频分箱 & 组员2 & 训练折拟合、逆变换或字段映射测试\\
探索与相关 & Pearson、Spearman、Kendall、分布和关联可视化 & 组员2 & 假设条件、显著性/效应量和非因果声明\\
聚类 & K-Means、DBSCAN、层次聚类、谱聚类 & 组员2 & 合成样例、簇稳定性和业务解释\\
插值与降维 & Lagrange、三次样条、Kriging；PCA、t-SNE、UMAP & 组员2 & 留点误差、解释方差、可视化边界\\
力学与物理 & Newton、弹簧阻尼、刚体平衡、摩擦、重力、Ohm/Kirchhoff、Bernoulli/质量守恒、理想气体 & 队长 & 受力/系统图、量纲、守恒/回代检查\\
ODE 模型 & Logistic、捕食—被捕食、Newton 冷却、药物代谢、SIR/SIS & 队长 & 初值、参数、解析/小样例和积分误差\\
PDE 模型 & 热、波、扩散、多孔流 & 队长 & 初边值、离散格式、稳定条件和网格收敛\\
数值方法 & Euler、RK4、FDM、FEM & 队长 & 同题对照、误差阶或收敛证据\\
规划 & LP、NLP、QP、整数、MILP、0--1、单/多目标、确定/随机模型 & 组员1 & 小规模枚举、界/间隙、约束映射\\
智能优化 & GA、PSO、SA、ACO、DE、贪心、NSGA-II、MOPSO、WOA、SSA、MFO & 组员1 & 同预算基线、20 种子、可行率和分布\\
图与网络 & Dijkstra、Floyd、SPFA、Prim、Kruskal、最大流、最小费用流、路径规划、二分图匹配、节点重要性 & 组员1 & 标准图测试、复杂度和负边/容量边界\\
仿真与敏感性 & Monte Carlo、情景仿真、参数扰动与策略切换 & 组员1 & 采样依据、重复次数、置信区间和稳定区间\\
主观评价 & AHP、模糊综合评价、FAHP & 组员2 & 判断矩阵、一致性、隶属函数与权重扰动\\
客观评价 & 灰色关联、PCA、因子分析、熵权、CRITIC、变异系数 & 组员2 & 正向化、尺度、权重来源和秩稳定性\\
综合评价 & TOPSIS、VIKOR、灰色综合评价、物元可拓、熵权--AHP、CRITIC--TOPSIS & 组员2 & 方法语义、基线排序和扰动对照\\
回归 & 多元线性、Ridge、Lasso、SVR、RF、XGBoost、CatBoost & 组员2 & pipeline、交叉验证、残差与外推边界\\
分类 & Logistic、朴素 Bayes、SVM、决策树、RF、AdaBoost、CNN & 组员2 & 混淆矩阵、阈值、类别不平衡与校准\\
时间序列 & ARIMA、SARIMA、指数平滑、VAR、GM(1,1)、灰色--Markov、Prophet、XGBoost、LightGBM、BP、RNN、LSTM、TCN、GRU & 组员2 & 持续/季节基线、滚动验证和预测区间\\
验证与误差 & MAE、MSE、RMSE、MAPE、K 折、留出、训练测试、分组/滚动验证 & 三人交叉 & 划分先于预处理；指标与任务损失一致\\
工程实现 & MATLAB、Python、配置、日志、随机种子、图表和结果导出 & 三人交叉 & 一键入口、依赖文件、单元测试和结果追溯\\
\bottomrule
\end{longtable}

\section{证据文件说明}

全量 OCR 页面记录、论文汇总、证据台账、模型链和语料统计属于临时分析证据，
保存在工作区 \texttt{.cumcm-work/} 下；正式成果不依赖运行时缓存才能编译。
台账中的页面文本可用于定位，不替代原 PDF。任何涉及公式、表格、小字号图注
或低置信页的结论，须回到原始页面局部核验后才能加入模板或 Skill。

\end{document}
